\documentclass[11pt]{article}
\usepackage{amsmath,amssymb,amsthm,mathtools}
\usepackage[margin=1in]{geometry}
\usepackage{enumitem}
\usepackage{hyperref}

\newtheorem{theorem}{Theorem}
\newtheorem{proposition}{Proposition}
\newtheorem{lemma}{Lemma}
\newtheorem{definition}{Definition}
\newtheorem{warning}{Warning}
\newtheorem{gate}{Gate}

\DeclareMathOperator{Ran}{Ran}
\DeclareMathOperator{Ker}{Ker}
\DeclareMathOperator{Cap}{Cap}
\DeclareMathOperator{Tr}{Tr}

\title{Step 115: Zeta-Factorization Attempt for the Boundary Block and the Source-Absorption Fallback}
\author{RATCHET RH Membrane Program}
\date{}

\begin{document}
\maketitle

\section{Purpose}
Steps 113--114 reduced coefficient visibility for the RH source route to the question whether the shifted Sonin/prolate boundary-packet sector lands in Burnol's co-Poisson complement.  In notation,
\[
  \mathcal B
  =\overline{\operatorname{span}}
  \{\operatorname{Ran}|P_\infty\tau_\ell(I-P_\infty)|:\ell\in\mathcal L_S\setminus\{0\}\}
\]
and Burnol's theorem gives
\[
  P_a=Y_a^\perp,
\]
where \(Y_a\) is the closed span of the zeta-zero evaluator vectors in the Sonine space \(L_a\).  The inclusion test is therefore
\[
  \mathcal B_a\subseteq P_a
  \quad\Longleftrightarrow\quad
  \Pi_{Y_a}\mathcal B_a=0.
\]
Step 115 attempts the strongest possible certificate for this inclusion: a zeta-factorization of one shifted boundary block.

\section{Boundary block}
Let \(J_a\) denote the legal transport from the semilocal/archimedean response geometry into Burnol's \(L_a\)-carrier.  For a log-shift mode \(\ell\), define
\[
  \mathfrak B_{\ell,a}=J_aP_\infty\tau_\ell(I-P_\infty).
\]
The Burnol co-Poisson inclusion gate is
\[
  \operatorname{Ran}\mathfrak B_{\ell,a}\subseteq P_a.
\]
Since \(P_a=Y_a^\perp\), this is equivalent to zero-evaluator annihilation:
\[
  \langle \mathfrak B_{\ell,a}u,Y^a_{\rho,k}\rangle=0
  \quad
  \text{for all legal }u,\rho,0\le k<m_\rho.
\]
Equivalently,
\[
  M(\mathfrak B_{\ell,a}u)^{(k)}(\rho)=0.
\]

\section{The zeta-factorization certificate}
Burnol's co-Poisson atoms have Mellin transforms containing a zeta factor.  The desired strong certificate is therefore:
\[
  M(\mathfrak B_{\ell,a}u)(s)=\zeta(s)\alpha_{\ell,u}(s),
\]
with \(\alpha_{\ell,u}\) legal in the appropriate Hardy/Sonine class and with all endpoint, pole, support, and trivial-zero records declared.

\begin{gate}[Zeta-factorization gate]
The shifted boundary block \(\mathfrak B_{\ell,a}\) passes the zeta-factorization gate if there exists a lawful linear map
\[
  A_\ell:u\mapsto \alpha_{\ell,u}
\]
such that
\[
  M(\mathfrak B_{\ell,a}u)=\zeta\,A_\ell u
\]
on the declared domain, with no target-zero information used to select \(A_\ell\).
\end{gate}

\begin{proposition}[Factorization implies inclusion]
If the zeta-factorization gate passes, then
\[
  \operatorname{Ran}\mathfrak B_{\ell,a}\subseteq P_a.
\]
\end{proposition}

\begin{proof}
For every nontrivial zero \(\rho\) of multiplicity \(m_\rho\), the factorization gives
\[
  M(\mathfrak B_{\ell,a}u)^{(k)}(\rho)=0,
  \qquad 0\le k<m_\rho,
\]
because \(\zeta\) has a zero of order \(m_\rho\) at \(\rho\).  Thus every zero-evaluator vector annihilates the range of \(\mathfrak B_{\ell,a}\), so the range is contained in \(Y_a^\perp=P_a\).
\end{proof}

\section{What the log-shift alone can and cannot do}
A log-shift does not itself create a zeta factor.  In Mellin/Fourier variables it is a nonvanishing exponential multiplier.  Schematically,
\[
  M(\tau_\ell f)(s)=e^{\ell\cdot\kappa(s)}M(f)(s),
\]
where the multiplier is nonzero at the nontrivial zeros.  Thus the shift cannot be the source of universal vanishing at \(\zeta\)-zeros.  If factorization holds, the vanishing must come from the Sonin/prolate projection geometry, co-Poisson synthesis, or an additional declared source/quotient record.

\begin{warning}[No automatic zeta divisor]
The implication
\[
  \text{log-shifted Sonin boundary packet}
  \Longrightarrow
  \zeta(s)\text{-divisible Mellin transform}
\]
is not automatic.  Treating the shift as if it produced a zeta factor is a smuggling error.
\end{warning}

\section{Residual if factorization fails}
If exact factorization is not proved, the unresolved visible-to-zero component is the positive residual
\[
  \Xi^{\rm BC}_{\ell,a}
  =\mathfrak B_{\ell,a}^{*}\Pi_{Y_a}\mathfrak B_{\ell,a}.
\]
This is a genuine adequacy residual: it measures the part of the boundary packet outside the declared Burnol/co-Poisson atom span.

\begin{theorem}[Boundary residual trichotomy]
For each legal shift mode \(\ell\), exactly one of the following proof statuses is available:
\begin{enumerate}[label=(\roman*)]
\item \textbf{Exact co-Poisson inclusion:}
\[
  \Xi^{\rm BC}_{\ell,a}=0.
\]
This follows, for example, from a lawful zeta-factorization.
\item \textbf{Defect-paid inclusion:}
\[
  \Xi^{\rm BC}_{\ell,a}\preceq E_{\ell,a},
\]
where the defect is declared and vanishes in a fixed/exhaustive ledger.
\item \textbf{Source absorption:}
\[
  \Xi^{\rm BC}_{\ell,a}\preceq F_n+E_{\ell,a,n},
\]
where \(F_n\) is an upstream-visible Hecke/Dirichlet source frame with a lower-frame record and the defects vanish.
\end{enumerate}
Without one of these records, the boundary-to-co-Poisson claim is support-only.
\end{theorem}

\section{Source-absorption fallback}
Let \(\mathcal R_{\rm BC}\) be the closure of the range of \((\Xi^{\rm BC})^{1/2}\).  Let
\[
  F_n=\sum_{\omega\in\mathcal X_n}\lambda_{\omega,n}Q_\omega^*\Theta_\omega^{-1}Q_\omega.
\]
A residual-specific source frame is accepted if
\[
  \Pi_{\mathcal R_{\rm BC}}F_n\Pi_{\mathcal R_{\rm BC}}
  \succeq
  \Lambda_n
  \Pi_{\mathcal R_{\rm BC}}(\Theta_0^-)^{-1}\Pi_{\mathcal R_{\rm BC}}
  -R_n^{\rm BC},
  \qquad \Lambda_n\to\infty.
\]
Then the residual is paid by the source ladder, provided the residual-sector tails vanish in the fixed/exhaustive sense.

\section{Decision of Step 115}
The zeta-factorization route is the cleanest possible route, but it is not earned by the raw shift block.  The shift mode supplies a nonvanishing Mellin multiplier, not a zeta divisor.  Therefore, the live proof obligation becomes:
\[
  \Pi_{Y_a}\mathfrak B_{\ell,a}=0
  \quad\text{by a genuine Sonin/co-Poisson theorem,}
\]
or else:
\[
  \Xi^{\rm BC}_{\ell,a}
  \text{ is absorbed by the Hecke/Dirichlet source ladder.}
\]

\section{Nonclaim boundary}
This step does not prove RH.  It does not prove \(\mathcal B\subseteq P_a\).  It does not prove that shifted Sonin boundary packets are co-Poisson.  It proves the exact gate: factorization implies inclusion, and failure to prove factorization produces a positive adequacy residual that must be source-absorbed.

\end{document}
